% Options for packages loaded elsewhere
\PassOptionsToPackage{unicode}{hyperref}
\PassOptionsToPackage{hyphens}{url}
\PassOptionsToPackage{dvipsnames,svgnames,x11names}{xcolor}
%
\documentclass[
  11pt,
  a4paper,
]{article}
\usepackage{amsmath,amssymb}
\usepackage{iftex}
\ifPDFTeX
  \usepackage[T1]{fontenc}
  \usepackage[utf8]{inputenc}
  \usepackage{textcomp} % provide euro and other symbols
\else % if luatex or xetex
  \usepackage{unicode-math} % this also loads fontspec
  \defaultfontfeatures{Scale=MatchLowercase}
  \defaultfontfeatures[\rmfamily]{Ligatures=TeX,Scale=1}
\fi
\usepackage{lmodern}
\ifPDFTeX\else
  % xetex/luatex font selection
  \setmainfont[]{Latin Modern Roman}
  \setmathfont[]{Latin Modern Math}
\fi
% Use upquote if available, for straight quotes in verbatim environments
\IfFileExists{upquote.sty}{\usepackage{upquote}}{}
\IfFileExists{microtype.sty}{% use microtype if available
  \usepackage[]{microtype}
  \UseMicrotypeSet[protrusion]{basicmath} % disable protrusion for tt fonts
}{}
\makeatletter
\@ifundefined{KOMAClassName}{% if non-KOMA class
  \IfFileExists{parskip.sty}{%
    \usepackage{parskip}
  }{% else
    \setlength{\parindent}{0pt}
    \setlength{\parskip}{6pt plus 2pt minus 1pt}}
}{% if KOMA class
  \KOMAoptions{parskip=half}}
\makeatother
\usepackage{xcolor}
\usepackage[top=25mm,bottom=24mm,left=27mm,right=27mm]{geometry}
\setlength{\emergencystretch}{3em} % prevent overfull lines
\providecommand{\tightlist}{%
  \setlength{\itemsep}{0pt}\setlength{\parskip}{0pt}}
\setcounter{secnumdepth}{-\maxdimen} % remove section numbering
% Mathematics-focused preprint layout: monochrome, Latin Modern OpenType text and math.
% Latin Modern fonts are supplied by a standard TeX installation, not embedded in this source archive.
\usepackage{microtype}
\usepackage{setspace}
\usepackage{titlesec}
\usepackage{fancyhdr}
\usepackage{needspace}
\setstretch{1.13}
\setlength{\parindent}{1.15em}
\setlength{\parskip}{3.25pt plus 1.0pt minus .6pt}
\titlespacing*{\section}{0pt}{17pt plus 4pt minus 2pt}{7pt plus 2pt}
\titlespacing*{\subsection}{0pt}{13pt plus 3pt minus 1pt}{5pt plus 1pt}
\titlespacing*{\subsubsection}{0pt}{10pt plus 2pt minus 1pt}{3pt plus 1pt}
\titleformat{\section}{\normalfont\bfseries\large}{\thesection}{0.7em}{}
\titleformat{\subsection}{\normalfont\bfseries\normalsize}{\thesubsection}{0.6em}{}
\titleformat{\subsubsection}{\normalfont\bfseries\normalsize}{\thesubsubsection}{0.6em}{}
\AtBeginDocument{%
  \setlength{\abovedisplayskip}{12pt plus 3pt minus 2pt}%
  \setlength{\belowdisplayskip}{12pt plus 3pt minus 2pt}%
  \setlength{\abovedisplayshortskip}{9pt plus 2pt minus 1pt}%
  \setlength{\belowdisplayshortskip}{10pt plus 2pt minus 1pt}%
  \pagestyle{plain}%
}
\setlength{\emergencystretch}{4em}
\setcounter{secnumdepth}{0}
\ifLuaTeX
  \usepackage{selnolig}  % disable illegal ligatures
\fi
\usepackage{bookmark}
\IfFileExists{xurl.sty}{\usepackage{xurl}}{} % add URL line breaks if available
\urlstyle{same}

\usepackage{longtable,booktabs,array,calc}
\usepackage{fvextra}
\RecustomVerbatimEnvironment{verbatim}{Verbatim}{breaklines=true,breakanywhere=true,fontsize=\footnotesize}
\usepackage{upquote}

\hypersetup{
  pdftitle={Symmetry constraints on elliptic-curve scalar observables},
  pdfauthor={Louis Desbiens (human initiator and curator)},
  pdfsubject={AI-developed mathematical research artifact; independently unverified candidate results},
  colorlinks=true,
  linkcolor={black},
  filecolor={black},
  citecolor={black},
  urlcolor={black},
  pdfcreator={LaTeX via pandoc}}

\title{Symmetry constraints on elliptic-curve scalar observables}
\usepackage{etoolbox}
\makeatletter
\providecommand{\subtitle}[1]{%
  \apptocmd{\@title}{\par {\large #1 \par}}{}{}
}
\makeatother
\subtitle{Parity, normalized division polynomials, collision quotients, and Fourier null models}
\author{}
\date{Public mathematical research artifact \textbar{} 31 October 2026}
\begin{document}
\maketitle

\textbf{Human initiator and curator:} Louis Desbiens\\
\textbf{AI mathematical research contributor:} OpenAI's ChatGPT\\
\textbf{Document classification:} Public AI-developed mathematical
research artifact -- independently unverified candidate results\\
\textbf{Institutional affiliation:} None claimed\\
\textbf{Manuscript license:} Creative Commons Attribution 4.0 International (CC BY 4.0)

\textbf{Status and reading guide.} This document records a
human-initiated, substantially AI-developed mathematical investigation.
The human curator directed the objectives but does not independently
certify its advanced mathematical arguments. The novel claims and
supplied proofs are candidate results awaiting independent specialist
verification; public availability does not certify correctness,
completeness or originality. This artifact develops scalar-parity
symmetry restrictions, division-polynomial normalization identities, and
a deliberately specified random-sign Fourier model. Elementary pairing
and counting arguments are included, while normalization conventions,
covariance calculations and the probabilistic limit should be checked
independently. The surrogate model is not a measured property of
deterministic elliptic-curve predicates, and no ECDLP hardness theorem
or cryptanalytic speedup is proposed.

\section{Abstract}\label{abstract}

This manuscript investigates the constraints imposed by exact symmetries
on public functions of elliptic-curve points when the corresponding
scalar is unknown. In a cyclic group of odd prime order, a
negation-invariant feature has one-half accuracy for canonical scalar
parity under a symmetric distribution; a pairing argument is supplied.
The treatment contrasts this obstruction with orientation-sensitive
contact quantities and presents a proposed normalization identity for an
intrinsic division-polynomial point function, distinct from a normalized
elliptic divisibility sequence. The sixfold complex-multiplication
quotient gives collision identities, but the idealized birthday gain is
only a constant-factor effect. A coefficient-counting argument limits
bounded-digit cyclotomic scalar representations. The manuscript also
derives Fourier symmetry identities and studies a deliberately defined
random-sign surrogate: its covariance is calculated and Gaussian
marginal behavior is argued using characteristic functions. Parseval
normalization, signed-orbit selection, and the limitations of
Fourier-envelope deductions are treated separately. The candidate
conclusions concern mathematical identities and the stated surrogate
model, not benchmarks, deterministic parity leakage, or a general ECDLP
lower bound.

\textbf{Keywords:} elliptic discrete logarithm; scalar parity; division
polynomial; complex multiplication; collision search; Fourier transform;
symmetry.

\section{1. Questions and mathematical
setting}\label{questions-and-mathematical-setting}

Let \(E/\mathbf F_q\) be a nonsingular short Weierstrass curve in
characteristic greater than three. Let \(G\) generate a subgroup of odd
prime order \(n>3\), and write \(P_d=[d]G\). A scalar means its residue
class modulo \(n\) unless a canonical representative \(0\leq d<n\) is
explicitly specified. Coordinate arithmetic takes place in
\(\mathbf F_q\), not in the scalar field \(\mathbf F_n\).

The exploratory investigation used numerical evaluations to test and
refine the initial intuition about coordinate observables. This note
isolates the exact identities and symmetry constraints underlying those
questions. It makes no quantitative claim about a historical classifier,
timing measurement, or numerical sample. The distinction is important:
an algebraic obstruction, a deliberately specified random model, and a
measured property of an actual deterministic predicate are different
statements.

The standard theory of elliptic arithmetic, division polynomials, and
collision algorithms is described in {[}1{]}. The relationship between
scalar parity, normalized elliptic divisibility sequences, and ECDLP is
prior work {[}2{]}. The role of efficient endomorphisms in known-scalar
multiplication is likewise established {[}3{]}. We do not claim those
general mechanisms as new.

\section{2. Scalar division and an exact parity
oracle}\label{scalar-division-and-an-exact-parity-oracle}

\textbf{Proposition 2.1 (scalar layers).} Let \(m\geq1\) be coprime to
\(n\). For \(0\leq k<n\), the representative \(x\) with \(0\leq x<n\) of
\(km^{-1}\pmod n\) satisfies

\[
x=\frac{k+\ell n}{m},\qquad 0\leq\ell<m,\qquad
k+\ell n\equiv0\pmod m. \qquad (2.1)
\]

\textbf{Proof.} The integer \(mx-k\) is a multiple of \(n\). It lies
strictly above \(-n\) and strictly below \(mn\), so it is \(\ell n\)
with \(0\leq\ell<m\). The congruence determines \(\ell\) uniquely since
\(n\) is invertible modulo \(m\). Conversely, that choice makes the
quotient integral and between zero and \(n-1\). \(\square\)

In particular, for \(m=2\), even \(k\) gives \(x=k/2<n/2\), whereas odd
\(k\) gives \(x=(k+n)/2>n/2\). This is a statement about the hidden
integer representative. The easily computed point \([m^{-1}]P_k\) does
not itself identify which scalar interval contains \(x\).

\textbf{Proposition 2.2 (an exact oracle suffices).} An oracle returning
canonical scalar parity for every point of \(\langle G\rangle\)
determines the complete scalar using at most \(\lceil\log_2 n\rceil\)
oracle calls and known group operations.

\textbf{Proof.} Starting from \(Q=[d]G\), let \(b\in\{0,1\}\) be the
returned parity and set \(Q'=[2^{-1}](Q-[b]G)\). Since \(d-b\) is even
as an integer and \(0\leq(d-b)/2<n\), the canonical scalar of \(Q'\) is
exactly \((d-b)/2\), with no modular wrap. Repeat until the point is
\(O\). The recorded bits, in least-significant-first order, recover
\(d\). This is the standard parity-oracle reduction, also given in {[}2,
Proposition 2{]}. \(\square\)

A bit of an ordinate is not that oracle. Its usefulness requires a
separate relation to the canonical scalar label.

\section{3. Negation-invariant features cannot predict
parity}\label{negation-invariant-features-cannot-predict-parity}

\textbf{Proposed Theorem 3.1.} Let \(L(P)=L(-P)\) on nonidentity
subgroup points. Under the uniform distribution on nonzero scalars,
every deterministic predictor depending only on \(L(P_d)\) has parity
accuracy exactly \(1/2\). The same holds in expectation for a randomized
predictor whose auxiliary randomness is independent of the hidden
scalar. More generally, the conclusion holds for any distribution
assigning equal mass to \(d\) and \(n-d\).

\textbf{Proof.} The involution \(d\mapsto n-d\) pairs all nonzero
scalars without fixed points. Because \(n\) is odd, the canonical
representatives in a pair have opposite parity. Their features are
identical, so a deterministic prediction is correct on exactly one
member. For a randomized predictor, condition on its random choices or
pair its two equal prediction probabilities. Averaging proves the
statement. \(\square\)

The theorem applies to any vector of functions of \(x(P)\) or
\(y(P)^2\), not only to one scalar feature. Undefined values must be
excluded symmetrically. It does not apply to a nonsymmetric interval
prior or to features containing orientation information. In particular,
adding a known offset to a point can destroy the symmetry of a feature,
but that fact alone does not imply predictive advantage.

\section{4. Contact norms and an orientation-sensitive
quantity}\label{contact-norms-and-an-orientation-sensitive-quantity}

Write \(E:y^2=C(x)\) and let \(A,B\) be fixed rational functions of
\(x\). For \(h=A-By\) and \(\bar h=A+By\),

\[
N=h\bar h=A^2-CB^2
\]

is invariant under negation. Wherever the denominators are nonzero,
define

\[
J_h=\frac{d}{dx}\log\frac{h}{\bar h}
=\frac{2C(A'B-AB')-ABC'}{y(A^2-CB^2)}. \qquad (4.1)
\]

\textbf{Derivation.} Since \(y'=C'/(2y)\), expand \(h'\bar h-h\bar h'\).
The cross terms leave \(2y(A'B-AB')-ABC'/y\). Dividing by \(h\bar h\)
yields (4.1). The numerator depends only on \(x\), while the denominator
contains one factor of \(y\), so \(J_h(-P)=-J_h(P)\).

The fixed-function qualification matters: if \(A,B\) are themselves
selected by a point-dependent rule, the whole rule's transformation
under negation must be checked. For fixed \(A,B\), the norm is subject
to Theorem 3.1 and the oriented quantity is not. Neither statement
supplies a correlation theorem for \(J_h\).

\section{5. An intrinsic normalized division-polynomial
function}\label{an-intrinsic-normalized-division-polynomial-function}

Assume in this section that \(\gcd(n,q-1)=1\). Let \(\psi_m\) be the
standard division polynomials, with \(\psi_1=1\), \(\psi_2=2y\), and
composition identity

\[
\psi_{ab}(P)=\psi_a([b]P)\,\psi_b(P)^{a^2}. \qquad (5.1)
\]

For nonzero \(P\in\langle G\rangle\), define \(L=q-1\), choose \(u\)
with \(u n^2\equiv1\pmod L\), and set

\[
\Theta(P)=\left(\frac{\psi_L(P)}{\psi_{L+n}(P)}\right)^u. \qquad (5.2)
\]

Both denominator and numerator are nonzero: their indices are not
divisible by the prime order \(n\). Exponentiation by \(u\) is the
unique inverse to the \(n^2\)-power map of \(\mathbf F_q^*\). Thus (5.2)
is a function of the point, with known indices independent of its
unknown scalar.

\textbf{Proposed Theorem 5.1 (normalization identity).} For every
integer \(e\) not divisible by \(n\),

\[
\Theta([e]P)=\Theta(P)^{e^2}\psi_e(P). \qquad (5.3)
\]

\textbf{Proof.} First take positive \(e\); negative indices follow from
\(\psi_{-e}=-\psi_e\) and the same composition identity. Put
\(U=\psi_e(P)\). Since \([L+n]P=[L]P\), applying (5.1) in both orders
gives

\[
\frac{\psi_L([e]P)}{\psi_{L+n}([e]P)}
=\left(\frac{\psi_L(P)}{\psi_{L+n}(P)}\right)^{e^2}
 U^{(L+n)^2-L^2}.
\]

The last exponent is congruent to \(n^2\) modulo \(L\). Taking the
unique \(n^2\)th root proves (5.3). All cancellations are legitimate:
\(n\) divides neither \(e\) nor \(L\), so \(\psi_e(P)\) and
\(\psi_e([L]P)\) are nonzero. \(\square\)

To distinguish the conventions explicitly, the normalized sequence
attached to a \emph{fixed} base point is

\[
\widetilde W_P(e)=\frac{\Theta([e]P)}{\Theta(P)}
=\Theta(P)^{e^2-1}\psi_e(P),\qquad \widetilde W_P(1)=1. \qquad (5.4)
\]

The exponent in the intrinsic point-function identity is \(e^2\); the
exponent in the normalized sequence is \(e^2-1\). The distinction is
forced already by substituting \(e=1\). With \(\Theta(O)=0\), (5.4) is
periodic in \(e\) modulo \(n\) and has zero terms exactly at multiples
of \(n\). The specific factor \(\Theta(P)^{e^2-1}\) preserves the
normalized Ward recurrence: both sides acquire the common factor
\(\Theta(P)^{2(m^2+n_0^2)-2}\) when its recurrence indices are
\(m,n_0\). This is the normalization mechanism studied in {[}2, Section
3{]}; the derivation above fixes the point-function and sequence
conventions without relying on a formula inferred from numerical data.

Let \(\chi_q\) be the quadratic character. If \(\chi_q(\Theta(P))=-1\),
(5.3) gives

\[
\chi_q(\Theta([e]P))=(-1)^e\chi_q(\psi_e(P)). \qquad (5.5)
\]

The left side is a public point function, but the unknown-index term on
the right is not supplied by it. Knowing that term would provide scalar
parity; merely computing \(\Theta(Q)\) does not. The character condition
on \(\Theta(P)\) is explicit, rather than assumed for every generator.
The normalization and residue problems in {[}2{]} explain why this is a
substantive missing input.

\section{6. Sixfold symmetry and exact collision
equations}\label{sixfold-symmetry-and-exact-collision-equations}

Assume \(E:y^2=x^3+b\) with \(b\neq0\), \(q\equiv1\pmod3\), and that the
subgroup is preserved by \(\phi(x,y)=(\beta x,y)\), where \(\beta\) has
order three. Write \(\phi(G)=[\lambda]G\) and assume \(\lambda\) has
order three modulo \(n\). These hypotheses hold for the standard
secp256k1 group {[}4{]}. Put \(H=\{\pm1,\pm\lambda,\pm\lambda^2\}\).
Every nonzero scalar orbit has six elements, since multiplication by a
nontrivial unit of \(H\) fixes no element of \(\mathbf F_n^*\).

A collision state with public coefficients is \(R_i=[A_i]G+[B_i]Q\),
where \(Q=[d]G\). If \(R_i=[v]R_j\) for a known \(v\in H\), then

\[
(B_i-vB_j)d=vA_j-A_i\pmod n.
\]

When \(B_i-vB_j\neq0\), this determines

\[
d=(vA_j-A_i)(B_i-vB_j)^{-1}\pmod n. \qquad (6.1)
\]

This follows by substituting \(Q=[d]G\) and comparing coefficients of a
generator. If the denominator is zero, a consistent collision also has
zero numerator and provides no equation for \(d\). Every proposed answer
should be checked by \([d]G=Q\).

Canonicalizing a state by \(R\mapsto[v]R\) requires changing its
coefficients to \((vA,vB)\) at the same time. Otherwise the stored
relation between point and coefficients is false.

\subsection{6.1. The idealized birthday
gain}\label{the-idealized-birthday-gain}

For independent uniform samples from a set of \(M\) states, the
probability of no collision in the first \(s\) samples is

\[
\prod_{j=0}^{s-1}(1-j/M).
\]

For \(s=O(\sqrt M)\), expanding its logarithm gives
\(-s(s-1)/(2M)+O(s^3/M^2)\). Thus the collision scale is \(\sqrt M\).
Replacing \(n-1\) states by \((n-1)/6\) in this \emph{sampling model}
yields the factor \(\sqrt6\) in sample count. It does not establish a
measured running-time gain for a deterministic or Pollard-style quotient
walk. Walk mixing, fruitless cycles, restarts, field costs, and
coefficient bookkeeping require their own analysis {[}1{]}.

Efficient-endomorphism decompositions \(k\equiv a+\lambda b\pmod n\) can
accelerate a known scalar multiplication \([k]P=[a]P+[b]\phi(P)\)
{[}3{]}. Applying the same notation to an unknown \(d\) does not reveal
\(a,b\) or produce two independent smaller discrete logarithms.

\section{7. Bounded-digit representations need a coverage
proof}\label{bounded-digit-representations-need-a-coverage-proof}

\textbf{Proposition 7.1.} Let a field of odd characteristic contain an
element \(\gamma\) of even order \(D\), and let
\(A_j\in\{0,\ldots,L-1\}\). The number of distinct values represented by
\(\sum_{j=0}^{D-1}A_j\gamma^j\) is at most

\[
(2L-1)^{D/2}. \qquad (7.1)
\]

\textbf{Proof.} Since \(\gamma^{D/2}=-1\),

\[
\sum_{j=0}^{D-1}A_j\gamma^j
=\sum_{j=0}^{D/2-1}(A_j-A_{j+D/2})\gamma^j.
\]

There are at most \(2L-1\) possible integer differences at each
position. Reducing them into the field or identifying different vectors
can only lower the number of outputs. \(\square\)

For \(D=64\) and \(L=16\), the upper bound is \(31^{32}\), fewer than
\(2^{159}\), rather than \(16^{64}=2^{256}\). This rules out full
coverage of a roughly 256-bit scalar field by that particular digit box.
Cycling coefficients can still give a valid multiplication identity;
coverage does not follow from that identity.

\section{8. Ordinate predicates and Fourier
symmetry}\label{ordinate-predicates-and-fourier-symmetry}

Now let \(q=p\) be an odd prime, let \(E(\mathbf F_p)=\langle G\rangle\)
have prime order \(n\), and retain the order-three automorphism of
Section 6. For canonical ordinate representatives \(0\leq y<p\), three
possible predicates are

\[
f_{\mathrm{bit}}(P)=(-1)^y,\qquad
f_{\mathrm{half}}(P)=\begin{cases}1,&2y<p,\\-1,&2y>p,\end{cases}
\] \[
f_{\mathrm{quad}}(P)=\chi_p(y(P)),\qquad f(O)=0.
\]

There is no rational nonzero two-torsion, so \(y\neq0\) away from \(O\).
All three predicates are invariant under \(\phi\). The first two are odd
under negation because the representative of \(-y\) is \(p-y\); the
third is odd when \(p\equiv3\pmod4\). None of these identities is a
scalar-parity prediction.

For any real sequence satisfying \(f_0=0\), \(f_{\lambda d}=f_d\), and
\(f_{-d}=-f_d\), define

\[
S_k=\sum_{d=0}^{n-1}f_d e^{-2\pi i kd/n}.
\]

Changing variables gives

\[
S_0=0,\quad S_{\lambda k}=S_k,\quad
S_{-k}=-S_k=\overline{S_k}. \qquad (8.1)
\]

The spectrum is purely imaginary, and there are \((n-1)/6\) nonzero
frequency orbits of equal magnitude. Equality of frequency values within
an orbit is deterministic dependence, not evidence for independent
observations.

\section{9. A symmetry-preserving random-sign
model}\label{a-symmetry-preserving-random-sign-model}

Choose one representative \(d_j\) of each sixfold scalar orbit and
independent signs \(\varepsilon_j\in\{-1,1\}\) of equal probability.
Define \(f_{\lambda^a d_j}=\varepsilon_j\),
\(f_{-\lambda^a d_j}=-\varepsilon_j\) for \(a=0,1,2\), and \(f_0=0\).
This is a specified surrogate, not an assertion that an arithmetic
predicate is random.

\textbf{Proposition 9.1 (covariance).} For nonzero frequencies \(k,l\),

\[
\mathbb E S_k=0,\qquad
\mathbb E(S_k\overline{S_l})=
\begin{cases}
n,&l\in\{k,\lambda k,\lambda^2k\},\\
-n,&l\in\{-k,-\lambda k,-\lambda^2k\},\\
0,&\text{otherwise}.
\end{cases} \qquad (9.1)
\]

\textbf{Proof.} For nonzero scalar indices, the input covariance is one
when \(e=\lambda^a d\), minus one when \(e=-\lambda^a d\), and zero
otherwise. Substitute this into the double Fourier sum. For each \(a\),
the sums over \(d\neq0\) are \(n-1\) when the relevant frequency is zero
and \(-1\) otherwise. The two constant terms cancel, leaving

\[
n\sum_{a=0}^2\bigl(\mathbf1_{l\lambda^a=k}-\mathbf1_{l\lambda^a=-k}\bigr).
\]

The three multipliers are distinct and none is the negative of another,
which proves (9.1). \(\square\)

\textbf{Proposition 9.2 (Gaussian marginals of the surrogate).} Along
any sequence of such prime orders \(n\to\infty\), the real variable
\(S_k/(i\sqrt n)\) converges in distribution to a standard normal
variable for any choice of a nonzero frequency at each size.

\textbf{Proof.} The contribution of each scalar orbit is an imaginary
number times its independent sign. Thus
\(S_k/i=\sum_j a_j\varepsilon_j\) with real \(a_j\) and \(|a_j|\leq6\).
Proposition 9.1 gives \(\sum_j a_j^2=n\). For fixed real \(t\), the
characteristic function is

\[
\prod_j\cos(t a_j/\sqrt n).
\]

Using \(\log\cos u=-u^2/2+O(u^4)\) uniformly near zero, its logarithm is

\[
-\frac{t^2}{2}+O\!\left(\frac{t^4\sum_j a_j^4}{n^2}\right)
=-\frac{t^2}{2}+O(t^4/n),
\]

since \(\sum_j a_j^4\leq36n\). Convergence of characteristic functions
proves the assertion. \(\square\)

This is a theorem about the explicitly generated signs. It is not a
Gaussian-limit theorem for ordinate parity, threshold predicates, or the
quadratic ordinate character. Nor does zero covariance between different
frequency orbits make the finite-size Fourier coefficients independent.

\section{10. Parseval, sign selection, and envelope
consequences}\label{parseval-sign-selection-and-envelope-consequences}

For any deterministic sign sequence with \(f_0=0\) and \(|f_d|=1\) for
\(d\neq0\), Fourier orthogonality gives

\[
\sum_{k=0}^{n-1}|S_k|^2=n(n-1). \qquad (10.1)
\]

When \(S_0=0\), its nonzero-frequency average is exactly \(n\). If
\(S_k\) is imaginary and \(X_k=S_k/(2i\sqrt p)\), then

\[
\frac1n\sum_k X_k^2=\frac{n-1}{4p},\qquad
\frac1{n-1}\sum_{k\neq0}X_k^2=\frac{n}{4p}. \qquad (10.2)
\]

Thus a variance match is not an independent distribution test. Selecting
one member from each sixfold frequency orbit preserves all even moments,
but not the signed empirical law: choosing the positive imaginary member
of each pair makes the selected values nonnegative. A canonical signed
comparison uses the full spectrum or includes both signs for each orbit.

\textbf{Proposition 10.1 (consequences of an assumed envelope).} Suppose
\(|S_k|\leq B\) at every nonzero frequency and \(S_0=0\). A
Fourier-supported approximation capturing a fraction \(\eta\) of the
total energy must use at least \(\eta n(n-1)/B^2\) frequencies. Also,
for every fixed \(u\neq0,v\pmod n\),

\[
\left|\frac1{n-1}\sum_{d=1}^{n-1}(-1)^d f_{ud+v}\right|
\leq\frac{B(2+\log n)+1}{n-1}. \qquad (10.3)
\]

\textbf{Proof.} A set of \(s\) frequencies contains at most \(sB^2\)
energy; compare with (10.1). For the second assertion let \(a_d=(-1)^d\)
on \(0\leq d<n\). A finite geometric sum gives

\[
|\widehat a(k)|=\frac1{|\cos(\pi k/n)|}.
\]

Pair frequencies around \(n/2\) and use \(\sin x\geq2x/\pi\) for
\(0\leq x\leq\pi/2\). The resulting sum over half-integers gives

\[
\sum_{k=1}^{n-1}|\widehat a(k)|\leq n(2+\log n).
\]

The Fourier magnitudes of \(d\mapsto f_{ud+v}\) are a permutation of
those of \(f\). Fourier inversion bounds the full scalar correlation by
\(B(2+\log n)\). Removing its \(d=0\) term changes it by at most one.
Division by \(n-1\) proves (10.3). \(\square\)

The envelope is an explicit hypothesis here, not an assumed property of
every predicate. These conclusions concern sparse Fourier approximation
and direct correlation with one fixed affine query. They do not rule out
arbitrary adaptive combinations, other features, or non-Fourier
algorithms.

\section{11. Interpretation}\label{interpretation}

The useful distinction is between information preserved by a symmetry
and a method for extracting a hidden scalar. Negation-invariant features
cannot predict canonical parity under symmetric sampling.
Orientation-sensitive quantities escape that exact obstruction but
require a further argument to become useful. Normalized
division-polynomial point functions satisfy a precise identity, yet the
unknown-index residue remains a separate input. Sixfold quotient
collisions yield exact equations when coefficients are tracked, while
their birthday benefit remains a constant-factor statement in an
explicit model.

The Fourier null model demonstrates how Gaussian-looking marginal
behavior can follow from symmetry-compatible random signs. It does not
prescribe the law of a deterministic arithmetic observable. Numerical
evaluations can motivate a sharper arithmetic theorem, but they neither
replace that theorem nor establish a general security conclusion. The
scope of every restriction in this note is therefore specified by its
hypotheses rather than by a numerical observation.

\section{Research provenance, contributions, and verification
status}\label{research-provenance-contributions-and-verification-status}

\subsection{Origin and development of the
investigation}\label{origin-and-development-of-the-investigation}

This mathematical investigation was initiated and directed by Louis
Desbiens and developed through extensive interaction with OpenAI's
ChatGPT. Louis Desbiens identified the original research questions, set
the objectives, guided the investigation, assessed its progress at a
conceptual level, and selected the material for public dissemination.

ChatGPT contributed substantially to mathematical exploration, including
alternative formulations, derivation of algebraic identities,
formulation and examination of conjectures and proposed proofs,
identification of potential counterexamples and exceptional cases,
design and execution of symbolic and numerical checks, organization of
relevant literature, and drafting and revision of the text. These
contributions went well beyond linguistic or editorial assistance.

\subsection{Attribution and scientific
authorship}\label{attribution-and-scientific-authorship}

Louis Desbiens is identified as the \textbf{human initiator and
curator}, not as an independent mathematical authority certifying the
results. OpenAI's ChatGPT is identified as the \textbf{AI mathematical
research contributor}, not as an accountable conventional scholarly
coauthor. These designations describe the provenance without implying
that the human curator personally derived or independently verified all
the arguments.

\subsection{Mathematical verification
status}\label{mathematical-verification-status}

The candidate mathematical contributions have not undergone independent
specialist validation, formal proof verification, or an exhaustive
originality and priority review. In particular, independent examination
is needed of the scope of the parity pairing; division-polynomial
normalization and recurrence conventions; the random-sign covariance and
Gaussian limit; separating the assumed Fourier envelope from
deterministic-curve conclusions. Locally displayed proofs are proposed
arguments. Established results cited from public literature should be
assessed separately from new deductions; neither an AI-generated proof
nor a successful computational check is by itself an independently
certified theorem. Individual statements may be correct, partly correct,
incomplete, or incorrect.

\subsection{Responsibility and purpose of public
release}\label{responsibility-and-purpose-of-public-release}

Louis Desbiens assumes responsibility for the decision to release this
collection and for a good-faith account of its origin, contributors and
verification status. He does \textbf{not} claim to have personally
established or independently certified the correctness, completeness or
originality of the proposed mathematical results. The material is made
available for independent examination, criticism, reproducibility and
possible further development; publication does not constitute scientific
endorsement. This disclosure does not eliminate ordinary
responsibilities associated with distributing public material.

\subsection{Institutional independence, funding, and competing
interests}\label{institutional-independence-funding-and-competing-interests}

No institutional affiliation or endorsement is claimed. The
investigation was undertaken independently, without institutional
funding or resources. No external research funding was received. The
human curator reports no competing interests.

\subsection{Computational evidence and source
availability}\label{computational-evidence-and-source-availability}

Numerical evaluations helped develop the initial intuition. The public
repository contains the available scripts, archived results, and a
claim-by-claim verification coverage guide. These are scoped finite
arithmetic, symbolic, and numerical checks, not independent verification
of the general mathematical results. Historical exploratory calculations
for which independent source data are unavailable are not treated as
reproducible evidence.

\section{References}\label{references}

{[}1{]} S. D. Galbraith, \emph{Mathematics of Public Key Cryptography},
Cambridge University Press, 2012. Public author version and errata:
\url{https://www.math.auckland.ac.nz/}\textasciitilde sgal018/crypto-book/crypto-book.html.

{[}2{]} K. E. Lauter and K. E. Stange, \emph{The elliptic curve discrete
logarithm problem and equivalent hard problems for elliptic divisibility
sequences}, arXiv:0803.0728v2, 2008; published in \emph{Selected Areas
in Cryptography}, LNCS 5381, 2009, pp.~309--327. DOI:
10.1007/978-3-642-04159-4\_20. Public version:
\url{https://arxiv.org/abs/0803.0728v2}. The composition argument and the
distinction between a point function and a normalized sequence are used
explicitly in the text.

{[}3{]} R. P. Gallant, R. J. Lambert, and S. A. Vanstone, ``Faster point
multiplication on elliptic curves with efficient endomorphisms,''
\emph{CRYPTO 2001}, LNCS 2139, 190--200, 2001. DOI:
10.1007/3-540-44647-8\_11. Public conference paper:
\url{https://www.iacr.org/archive/crypto2001/21390189.pdf}. A public
exposition of the endomorphism method is also contained in {[}1{]}.

{[}4{]} Certicom Research, \emph{SEC 2: Recommended Elliptic Curve
Domain Parameters}, version 2.0, 27 January 2010, Section 2.4.1.
\url{https://www.secg.org/sec2-v2.pdf}.
\end{document}
